\documentclass[10pt,a4paper]{article}

\usepackage[margin=2.2cm]{geometry}
\usepackage{array}
\usepackage{longtable}
\usepackage{parskip}
\usepackage{hyperref}

\title{GF2 Software Project: First Interim Report}
\author{
Team Number: \textit{Insert team number here}\\
\textit{Insert Name 1 and CRSID}\\
\textit{Insert Name 2 and CRSID}\\
\textit{Insert Name 3 and CRSID}
}
\date{24 May 2026}

\begin{document}

\maketitle

\section{Introduction and General Approach}

The aim of this project is to develop a Python logic simulator for
combinational and clocked logic circuits. The simulator will work in two main
phases. First, it reads a text definition file describing the devices in the
circuit, the connections between them, and the initial monitor points. Second,
it allows the user to run the network, change switch states, and inspect the
monitored output signals.

This report specifies the logic description language used in the definition
file. The language is designed to be readable, simple to parse, and easy to
check for errors. The syntax is deliberately fairly general: it defines the
overall structure of the file, while detailed device-specific checks are left
to semantic analysis. For example, the grammar can recognise a pin name such as
\texttt{G1.3}, but semantic checking decides whether pin \texttt{3} actually
exists on device \texttt{G1}.

The definition file is divided into an optional import block, a device block, a
connection block, a monitor block, and a final \texttt{END}. The main design
choice is to keep the syntax compact while using clear semantic error messages
to explain invalid circuits.

\section{Teamwork Planning}

The team divided the first interim work into three parallel tasks.

\begin{itemize}
    \item \textbf{Person A: Language syntax and scanner planning.}
    This person is responsible for the EBNF grammar, reserved words,
    punctuation, comments, whitespace handling, and scanner token decisions.

    \item \textbf{Person B: Semantics and error handling.}
    This person is responsible for defining semantic constraints, identifying
    semantic errors, and specifying how syntax and semantic errors will be
    reported.

    \item \textbf{Person C: Examples, diagrams, and report assembly.}
    This person is responsible for writing example definition files, drawing
    circuit diagrams, assembling the report, and checking consistency between
    the examples, grammar, and semantic rules.
\end{itemize}

The team will review the language together before implementation begins. This
is important because the language affects the scanner, parser, test files, GUI
examples, and user guide. During implementation, each person will mainly work
on their own module, but tests should be reviewed by someone other than the
module author where possible.

\begin{center}
\begin{tabular}{|p{0.18\linewidth}|p{0.72\linewidth}|}
\hline
\textbf{Date} & \textbf{Planned work}\\
\hline
15--17 May & Agree language structure, draft EBNF, choose examples.\\
\hline
18--21 May & Finalise grammar, semantic rules, error handling, and diagrams.\\
\hline
22--24 May & Assemble, proofread, and submit first interim report.\\
\hline
25 May--4 June & Implement names, scanner, parser, GUI, and tests.\\
\hline
5 June & Submit second interim work and receive maintenance task.\\
\hline
\end{tabular}
\end{center}

\section{Logic Description Language Syntax}

\subsection{General File Structure}

A circuit definition consists of the following parts:

\begin{enumerate}
    \item an optional \texttt{IMPORT} block for custom device types;
    \item a \texttt{DEVICES} block declaring all devices;
    \item a \texttt{CONNECT} block declaring all connections;
    \item a \texttt{MONITOR} block declaring initial monitor points;
    \item a final \texttt{END}.
\end{enumerate}

The basic shape of a file is:

\begin{verbatim}
DEVICES:
  ...
DEVICES END;

CONNECT:
  ...
CONNECT END;

MONITOR:
  ...
MONITOR END;

END
\end{verbatim}

A connection statement has the form:

\begin{verbatim}
OutputSignal = InputSignal;
\end{verbatim}

This means that the output signal on the left is connected to the input signal
on the right.

\subsection{EBNF Grammar}

The syntax of the language is defined by the following EBNF grammar.

\begin{verbatim}
CircuitDefinition = [ ImportsBlock ] , DevicesBlock,
                    ConnectionsBlock, MonitorsBlock, "END" ;

ImportsBlock      = "IMPORT", ":", { ImportRule },
                    "IMPORT", "END", ";" ;

ImportRule        = CustomTypeName, "FROM", QuotedString, ";" ;

DevicesBlock      = "DEVICES", ":", { DeviceDeclaration },
                    "DEVICES", "END", ";" ;

DeviceDeclaration = DeviceName, "=", DeviceType, [ DeviceParameter ], ";" ;

DeviceType        = PrimitiveType | CustomTypeName ;

PrimitiveType     = "SWITCH" | "CLOCK" | "AND" | "OR" | "NAND" |
                    "NOR" | "XOR" | "NOT" | "DTYPE" ;

CustomTypeName    = Letter, { Letter | Digit } ;

DeviceParameter   = Number ;

ConnectionsBlock  = "CONNECT", ":", { ConnectionRule },
                    "CONNECT", "END", ";" ;

ConnectionRule    = OutputSignal, "=", InputSignal, ";" ;

OutputSignal      = DeviceName, [ "." , PinName ] ;

InputSignal       = DeviceName, "." , PinName ;

MonitorsBlock     = "MONITOR", ":", { MonitorRule },
                    "MONITOR", "END", ";" ;

MonitorRule       = OutputSignal, ";" ;

DeviceName        = Letter, { Letter | Digit } ;

PinName           = ( Letter | Digit ) , { Letter | Digit } ;

Number            = Digit, { Digit } ;

QuotedString      = '"', { Character }, '"' ;

Letter            = "A" | "B" | ... | "Z" | "a" | ... | "z" ;

Digit             = "0" | "1" | "2" | "3" | "4" |
                    "5" | "6" | "7" | "8" | "9" ;
\end{verbatim}

Here, \texttt{Character} means any printable character except a double quote or
newline. This is used only inside quoted import paths.

\subsection{Comments and Whitespace}

Comments are handled by the scanner and are therefore not part of the formal
grammar. A comment begins with \texttt{\#} and continues to the end of the
current line.

\begin{verbatim}
G1 = NAND 2; # two-input NAND gate
\end{verbatim}

Whitespace is free-format. Spaces, tabs, and line breaks may appear between
tokens without changing the meaning of the file. Whitespace may not appear
inside names, numbers, quoted strings, or punctuation symbols.

\subsection{Scanner Tokens}

The scanner will translate the input file into a stream of tokens for the
parser. The required token types are:

\begin{center}
\begin{tabular}{>{\ttfamily}ll}
KEYWORD & reserved word such as DEVICES or SWITCH\\
NAME & user-defined name or pin name\\
NUMBER & integer number\\
STRING & quoted import path\\
COLON & \texttt{:}\\
SEMICOLON & \texttt{;}\\
EQUALS & \texttt{=}\\
DOT & \texttt{.}\\
EOF & end of file\\
INVALID & invalid character or token\\
\end{tabular}
\end{center}

The reserved keywords are:

\begin{verbatim}
IMPORT FROM DEVICES CONNECT MONITOR END
SWITCH CLOCK AND OR NAND NOR XOR NOT DTYPE
\end{verbatim}

Reserved keywords may not be used as device names or custom type names.

\section{Semantic Constraints}

The grammar describes the shape of a valid file, but not every grammatically
valid file describes a valid circuit. Semantic analysis checks whether the
devices, parameters, pins, connections, monitor points, and imports are
meaningful.

\subsection{Import Constraints}

\begin{itemize}
    \item The \texttt{IMPORT} block is optional.
    \item A custom type must be imported before it is used as a
    \texttt{DeviceType}.
    \item A custom type name must not be the same as a primitive type or
    reserved keyword.
    \item Two imports must not define the same custom type name.
    \item The quoted import path must refer to a readable file.
    \item Recursive imports are not allowed.
    \item If custom device simulation is not implemented, imports will be
    rejected with a clear semantic error.
\end{itemize}

\subsection{Device Constraints}

\begin{itemize}
    \item Every device name must be unique.
    \item Every device type must be a primitive type or an imported custom
    type.
    \item \texttt{SWITCH} requires one parameter, either \texttt{0} or
    \texttt{1}, giving the initial output state.
    \item \texttt{CLOCK} requires one positive integer parameter giving the
    number of simulation cycles in half a period.
    \item \texttt{AND}, \texttt{OR}, \texttt{NAND}, and \texttt{NOR} require
    one integer parameter giving the number of inputs. This value must be
    between 1 and 16.
    \item \texttt{XOR} has two inputs and must not be given a parameter.
    \item \texttt{NOT} has one input and must not be given a parameter.
    \item \texttt{DTYPE} has fixed inputs and outputs and must not be given a
    parameter.
\end{itemize}

\subsection{Pin and Signal Constraints}

\begin{itemize}
    \item Single-output devices may be used as output signals by writing the
    device name alone.
    \item Gate inputs may be referred to using numbered pins such as
    \texttt{1}, \texttt{2}, or named pins such as \texttt{I1}, \texttt{I2},
    depending on the implementation convention chosen.
    \item \texttt{DTYPE} input pins are \texttt{DATA}, \texttt{CLK},
    \texttt{SET}, and \texttt{CLEAR}.
    \item \texttt{DTYPE} output pins are \texttt{Q} and \texttt{QBAR}.
    \item A pin name is only semantically valid if it exists on the referenced
    device.
    \item A signal used in the left side of a connection must be an output.
    \item A signal used in the right side of a connection must be an input.
    \item A monitor point must be an output signal.
\end{itemize}

\subsection{Connection Constraints}

\begin{itemize}
    \item Each connection connects exactly one output signal to one input
    signal.
    \item An input may be connected at most once.
    \item One output may drive more than one input.
    \item A connection must not refer to an undeclared device.
    \item A connection must not refer to a non-existent pin.
    \item All required inputs must be connected before the simulation can run.
\end{itemize}

\subsection{Monitor Constraints}

\begin{itemize}
    \item A monitor point must refer to an existing device.
    \item A monitor point must refer to an output, not an input.
    \item Duplicate monitor points will be rejected as semantic errors.
    \item An empty monitor block is syntactically valid; monitors can be added
    later through the user interface.
\end{itemize}

\section{Error Conditions and Error Handling}

\subsection{Syntax Errors}

Syntax errors occur when the token sequence does not match the EBNF grammar.
Examples include:

\begin{itemize}
    \item missing section keyword;
    \item missing colon after a section keyword;
    \item missing semicolon after a rule;
    \item missing equals sign in a device declaration;
    \item missing equals sign in a connection rule;
    \item missing block terminator such as \texttt{DEVICES END;};
    \item malformed quoted string;
    \item invalid character;
    \item unexpected end of file.
\end{itemize}

When a syntax error is found, the scanner will print the current line and a
caret marker showing the approximate error position. The parser will then print
a short message explaining what was expected.

\begin{verbatim}
DEVICES:
  G1 NAND 2;
     ^
*** Error: Expected '=' after device name.
\end{verbatim}

After reporting a syntax error, the parser will skip tokens until it reaches a
safe recovery point. Recovery points include semicolons, block-ending keywords,
and end of file. This should allow several errors in one file to be reported in
a single run.

\subsection{Semantic Errors}

Semantic errors occur when the syntax is valid but the circuit description is
not meaningful. The main semantic errors are:

\begin{itemize}
    \item duplicate device name;
    \item reserved keyword used as a custom type or device name;
    \item custom type used without import;
    \item import file missing or invalid;
    \item unsupported primitive device;
    \item missing required parameter;
    \item unexpected parameter;
    \item invalid switch value;
    \item invalid clock period;
    \item invalid gate input count;
    \item unknown device in a connection or monitor;
    \item unknown pin;
    \item input used where an output is required;
    \item output used where an input is required;
    \item input connected more than once;
    \item monitor point refers to an input;
    \item required input left unconnected.
\end{itemize}

Semantic error messages will name the device, pin, or type causing the problem.

\begin{verbatim}
CONNECT:
  SW1 = G1.3;
        ^
*** Error: G1.3 is not valid because G1 has only two inputs.
\end{verbatim}

If any syntax or semantic error is found, parsing fails and the simulator will
not run. The parser may continue checking later statements to find additional
errors, but the network will only be executed if no errors are detected.

\section{Example Definition Files}

\subsection{Example 1: NAND Latch}

This example describes a cross-coupled NAND latch. Two switches provide the
external inputs and the outputs of the two NAND gates are monitored.

\begin{verbatim}
DEVICES:
  SW1 = SWITCH 0;
  SW2 = SWITCH 0;
  G1 = NAND 2;
  G2 = NAND 2;
DEVICES END;

CONNECT:
  SW1 = G1.1;
  SW2 = G2.2;
  G1 = G2.1;
  G2 = G1.2;
CONNECT END;

MONITOR:
  G1;
  G2;
MONITOR END;

END
\end{verbatim}

Diagram:

\begin{verbatim}
        +---------+                 +---------+
SW1 --->| G1 NAND |---- G1 -------->| G2.1    |
        | 1     2 |                 | NAND    |---- G2
G2  --->|         |                 | 2       |
        +---------+                 +---------+
                                           ^
                                           |
                                          SW2

Monitor points: G1 and G2
\end{verbatim}

\subsection{Example 2: Clocked D-Type Circuit}

This example uses a \texttt{DTYPE} device, three switches, and one clock. It
demonstrates named input pins and named output pins.

\begin{verbatim}
DEVICES:
  DATA = SWITCH 0;
  SETSW = SWITCH 0;
  CLRSW = SWITCH 0;
  CLK1 = CLOCK 2;
  D1 = DTYPE;
DEVICES END;

CONNECT:
  DATA = D1.DATA;
  CLK1 = D1.CLK;
  SETSW = D1.SET;
  CLRSW = D1.CLEAR;
CONNECT END;

MONITOR:
  D1.Q;
  D1.QBAR;
  CLK1;
MONITOR END;

END
\end{verbatim}

Diagram:

\begin{verbatim}
DATA  ----------------> DATA        Q -----> monitor D1.Q
CLK1  ----------------> CLK     +-------+
SETSW ----------------> SET     | DTYPE |
CLRSW ----------------> CLEAR   +-------+
                                  QBAR --> monitor D1.QBAR

CLK1 has half-period 2 and is also monitored.
\end{verbatim}

\section{Conclusion}

The proposed language gives a clear overall structure for logic definition
files while leaving device-specific validity checks to semantic analysis. This
keeps the grammar compact and allows the parser to report useful errors for
invalid devices, parameters, pins, connections, imports, and monitor points.
The next stage of the project will be to implement the names, scanner, parser,
and GUI modules using this specification.

\end{document}
